\[ %%%%%%%%%%%%%%%%%%%%%%%%%%%%%%%%%%%%%%%%%%%%%%%%%%%%%%%%%%%%%%%%%%%%%%%%%%%%%%%% \subsection{Dynamic Graph Representation of Physical Systems} \label{sec:dynamic_graph} %%%%%%%%%%%%%%%%%%%%%%%%%%%%%%%%%%%%%%%%%%%%%%%%%%%%%%%%%%%%%%%%%%%%%%%%%%%%%%%% A fundamental assumption of PHYSGEN is that complex physical systems can be represented as dynamic interaction graphs whose structure and states evolve over time. Unlike static graph neural networks, where the graph topology is predefined and remains unchanged during inference, PHYSGEN considers the graph itself as a dynamic representation of the underlying physical process. Let a physical system at time $t$ be represented by the graph: \begin{equation} G_t=(V_t,E_t,X_t,P_t), \end{equation} where: \begin{itemize} \item $V_t=\{v_1,v_2,\dots,v_N\}$ denotes the set of physical entities; \item $E_t$ represents interaction relationships between entities; \item $X_t$ contains dynamic physical state variables; \item $P_t$ represents physical parameters and structural properties. \end{itemize} Each graph node corresponds to a physical degree of freedom or computational entity. Depending on the application domain, nodes may represent spatial discretization points, particles, material components, biological units, or elements of engineered systems. The node state is defined as: \begin{equation} x_i^t= [ \rho_i^t, u_i^t, T_i^t, m_i, p_i, \ldots ], \end{equation} where the components may include physical quantities such as density, velocity, temperature, mass, pressure, material properties, or other domain-specific variables. The edge set describes physical interactions between system components: \begin{equation} E_t= \{e_{ij}^{t}|v_i,v_j\in V_t\}. \end{equation} Each edge is associated with an interaction feature: \begin{equation} e_{ij}^{t} = [ d_{ij}, r_{ij}, \omega_{ij}, c_{ij}, \ldots ], \end{equation} where possible components include spatial distance, relative position, interaction strength, connectivity information, or physical coupling coefficients. %%%%%%%%%%%%%%%%%%%%%%%%%%%%%%%%%%%%%%%%%%%%%%%%%%%%%%%%%%%%%%%%%%%%%%%%%%%%%%%% \subsubsection{Dynamic Topology Evolution} %%%%%%%%%%%%%%%%%%%%%%%%%%%%%%%%%%%%%%%%%%%%%%%%%%%%%%%%%%%%%%%%%%%%%%%%%%%%%%%% A distinguishing feature of PHYSGEN is that the graph topology is not assumed to be fixed. Instead, the interaction structure may evolve according to changes in the physical system. The graph transition can be expressed as: \begin{equation} E_{t+1} = \Gamma_{\theta}(E_t,X_t,P_t), \end{equation} where $\Gamma_{\theta}$ represents a learnable physics-guided topology evolution operator. This formulation allows PHYSGEN to represent systems where interactions change over time, including: \begin{itemize} \item moving particle systems; \item adaptive computational meshes; \item phase transitions; \item biological networks; \item evolving material structures. \end{itemize} Unlike conventional GNN approaches that operate on predetermined connectivity patterns, PHYSGEN allows physical interactions themselves to become part of the learned dynamical process. %%%%%%%%%%%%%%%%%%%%%%%%%%%%%%%%%%%%%%%%%%%%%%%%%%%%%%%%%%%%%%%%%%%%%%%%%%%%%%%% \subsubsection{Physics-Based Edge Modulation} %%%%%%%%%%%%%%%%%%%%%%%%%%%%%%%%%%%%%%%%%%%%%%%%%%%%%%%%%%%%%%%%%%%%%%%%%%%%%%%% In many physical systems, not all interactions contribute equally to system evolution. The influence of neighboring entities depends on physical conditions, conservation principles, and local state relationships. Therefore, PHYSGEN introduces a physics-guided edge modulation mechanism: \begin{equation} \tilde{e}_{ij}^{t} = \alpha_{ij}^{phys} e_{ij}^{t}, \end{equation} where $\alpha_{ij}^{phys}$ represents a physically constrained interaction coefficient. The modulation coefficient can be expressed as: \begin{equation} \alpha_{ij}^{phys} = f_{phys} ( x_i^t, x_j^t, e_{ij}^{t}, P_t ), \end{equation} where $f_{phys}$ combines learned representations with embedded physical knowledge. This mechanism allows the architecture to distinguish between physically relevant and insignificant interactions while preserving the flexibility of data-driven learning. %%%%%%%%%%%%%%%%%%%%%%%%%%%%%%%%%%%%%%%%%%%%%%%%%%%%%%%%%%%%%%%%%%%%%%%%%%%%%%%% \subsubsection{Graph-Based Approximation of Physical Dynamics} %%%%%%%%%%%%%%%%%%%%%%%%%%%%%%%%%%%%%%%%%%%%%%%%%%%%%%%%%%%%%%%%%%%%%%%%%%%%%%%% Many physical systems are described mathematically through continuous differential equations: \begin{equation} \frac{\partial u}{\partial t} = \mathcal{F}(u,\nabla u,\nabla^2u,\theta), \end{equation} where $\mathcal{F}$ represents the governing physical operator. After spatial discretization, such systems can be interpreted as interacting local states: \begin{equation} x_i^{t+1} = x_i^t + \Delta t \, \mathcal{F}_i ( x_i^t, \{x_j^t:j\in\mathcal{N}(i)\} ). \end{equation} PHYSGEN generalizes this formulation by learning a physics-guided graph evolution operator: \begin{equation} X_{t+1} = \Phi_{\theta}(G_t,X_t,P_t). \end{equation} Thus, the graph representation acts as an intermediate computational language between continuous physical equations and neural learning architectures. %%%%%%%%%%%%%%%%%%%%%%%%%%%%%%%%%%%%%%%%%%%%%%%%%%%%%%%%%%%%%%%%%%%%%%%%%%%%%%%% \subsubsection{Physical Invariants and Graph State Constraints} %%%%%%%%%%%%%%%%%%%%%%%%%%%%%%%%%%%%%%%%%%%%%%%%%%%%%%%%%%%%%%%%%%%%%%%%%%%%%%%% A critical requirement for reliable physical prediction is preservation of system invariants. Depending on the domain, these may include conservation of mass, momentum, energy, charge, probability, or other quantities. PHYSGEN introduces invariant-aware constraints through a physical consistency functional: \begin{equation} \mathcal{C}(X_t) = ||I(X_t)-I(X_{t+1})||, \end{equation} where $I(\cdot)$ represents a system invariant. The graph evolution process is therefore optimized under both predictive and physical objectives: \begin{equation} \min_{\theta} \mathcal{L} = \mathcal{L}_{prediction} + \lambda_c \mathcal{C}(X_t). \end{equation} However, unlike classical physics-informed approaches, these constraints are not the only mechanism of physical integration. In PHYSGEN, invariant preservation operates together with physics-guided message passing, latent dynamics regulation, and stability control. %%%%%%%%%%%%%%%%%%%%%%%%%%%%%%%%%%%%%%%%%%%%%%%%%%%%%%%%%%%%%%%%%%%%%%%%%%%%%%%% \subsubsection{Summary} %%%%%%%%%%%%%%%%%%%%%%%%%%%%%%%%%%%%%%%%%%%%%%%%%%%%%%%%%%%%%%%%%%%%%%%%%%%%%%%% The dynamic graph representation provides the mathematical foundation of PHYSGEN. By treating both node states and interaction structures as evolving physical quantities, the framework extends conventional graph neural networks toward adaptive representations of complex dynamical systems. This formulation enables PHYSGEN to model heterogeneous physical processes while preserving the relational structure, physical constraints, and temporal dependencies required for reliable long-horizon prediction. \] 
